\section{The invariant, its ideals, and matrix conventions}
\label{sec:structure}

The trace cone supplies three inputs to the construction: pointwise
comparison of Cuntz classes in \(Q\), a common lattice of ideals for
\(P\) and \(Q\), and positive elements that control finite collections
of rank estimates inside those ideals. We establish these inputs
before fixing the matrix identifications used in the approximations.
All later comparisons are obtained from the coordinate algebra
\(Q\); passage to a sequence algebra will require bounded witnesses.

\subsection{Ranks and realification}

For \(a\in(C\otimes\K)_+\), extend \(\tau\in T(C)\) to matrices by
the unnormalized matrix trace and put
\begin{equation}\label{st:rank}
 d_\tau(a)=\lim_{k\to\infty}\tau(a^{1/k})
 \qquad(0\leq a\leq1).
\end{equation}
Rescaling an element to be contractive does not change its rank.
The resulting function depends only on its Cuntz class. We write
\(d_\tau(x)\) also for the value at \(x\in\Cu(C)\).
The basic cut and integration identities are
\begin{equation}\label{st:rank-cuts}
 d_\tau(a)=\sup_{\eps>0}d_\tau((a-\eps)_+),
 \qquad
 \tau(a)=\int_0^{\|a\|}d_\tau((a-t)_+)\,dt.
\end{equation}
These identities use extended nonnegative arithmetic and apply to all
the weights in \eqref{intro:trace-topology}.

We use the following form of the trace and realification theorems.
The assertion about the functional cone is
\cite[Proposition~4.2, Remark~4.3, and Theorem~4.4]{ERS}; compactness
is \cite[Theorem~3.7]{ERS}. The realification assertions are
\cite[Proposition~3.1.1 and Theorems~3.2.1 and~5.1.2]{Robert}.

\begin{theorem}[Trace realification]\label{st:realification}
For every \(C^*\)-algebra \(C\), the trace cone with topology
\eqref{intro:trace-topology} is compact Hausdorff, and it is second
countable when \(C\) is separable. If \(C\) is exact, the assignment
\(\tau\mapsto d_\tau\) identifies this cone with the topological
cone of functionals on \(\Cu(C)\).

For every \(C^*\)-algebra \(C\),
\[
 \Cu(C\otimes\W)\cong\Cu(C)_{\R}.
\]
Realification is determined functorially by the topological functional
cone. Its realization as functions on that cone has pointwise order,
pointwise addition, and multiplication by positive real scalars.
It has the same functional cone as \(\Cu(C)\).

Consequently the given trace-cone isomorphism induces isomorphisms
\begin{equation}\label{st:transport}
 F:T(Q)\longrightarrow T(P),\qquad
 G:\Cu(P)\longrightarrow\Cu(Q),\qquad
 d_\tau(Gx)=d_{F\tau}(x).
\end{equation}
In particular, for \(x,y\in\Cu(Q)\),
\begin{equation}\label{st:pointwise-order}
 x\leq y
 \quad\Longleftrightarrow\quad
 d_\tau(x)\leq d_\tau(y)\quad\hbox{for every }\tau\in T(Q).
\end{equation}
Multiplication by every positive real scalar is an order automorphism
preserving increasing suprema and compact containment.
\end{theorem}

Here the function realization is Robert's cone
\(L(\operatorname{F}(\Cu(C)))\), where
\(\operatorname{F}(\Cu(C))\) denotes the functional cone. For a
functional cone \(T\), the cone \(L(T)\) consists of pointwise suprema
\(f=\sup_n f_n\), where the \(f_n:T\to[0,\infty]\) are additive,
positively homogeneous, lower semicontinuous, and zero at zero.
The approximating sequence satisfies, for some \(0<\eps_n<1\),
\[
 f_n\leq(1-\eps_n)f_{n+1},
 \qquad f_n\text{ is continuous at every }\tau\text{ with }f_{n+1}(\tau)<\infty.
\]
Continuity here is on all of \(T\) at the indicated point, with the
usual topology on \([0,\infty]\), rather than continuity only of a
restriction to the finite locus.
For the exact algebras considered below the realization is \(L(T(C))\).
This specified approximation condition is essential; arbitrary
lower-semicontinuous homogeneous functions are not being identified
with Cuntz classes. Composition with a cone homeomorphism preserves
all the displayed conditions, which is the functoriality used here. This gives \(G\) in
\eqref{st:transport}. We fix the normalized trace of \(\W\) when
identifying the trace cones before and after tensoring; stabilization
uses the unnormalized trace of \(\K\).
The equality of extended lower-semicontinuous quasitraces and traces
for exact algebras is included in the cited trace theorem. Its ancestry
is Haagerup's theorem for unital exact algebras
\cite[Theorem~5.11]{Haagerup} and the Blanchard--Kirchberg reduction
recalled in \cite[Remark~4.3]{ERS}. We use the extended nonunital
interface stated there. Thus the functional cone used here is exactly
the invariant of the problem.

We also record the tensor-factor properties
\begin{equation}\label{st:W-properties}
 \W\otimes\UHF\cong\W,
 \qquad
 \W\otimes\W\cong\W,
 \qquad
 \KK(\W,\W)=0.
\end{equation}
The first property follows from
\cite[Proposition~6.2]{Jacelon}; the remaining properties are recorded
in \cite[Section~2.1 and Corollary~6.2]{Nawata}.
They are properties of the tensor factor \(\W\), so their use places
no simplicity restriction on \(A,B,P\), or \(Q\).
The UHF absorption is used in the foldings below, and the
ideal-related use of the last equality is proved in
\Cref{un:KK-nullity}.

For completeness, the tensor-square and approximate-inner-flip
properties underlying realification have accessible proofs in these
same sources. Choose $\beta:\W\otimes\W\to\W$ by
\cite[Corollary~6.2(i)]{Nawata}. Conjugating the flip by $\beta$
gives a trace-preserving automorphism of $\W$, which is approximately
inner by \cite[Corollary~4.6]{Jacelon}. Pulling the implementing
unitaries back through the unitized $\beta$ gives approximate inner
flip of $\W\otimes\W$, as in \cite[Remark~5.2]{Jacelon}.
Thus the tensor-square isomorphism and the approximate inner flip
needed for realification follow from properties of the factor itself.

\subsection{Finite ideals, zero ideals, and open supports}

For \(\tau\in T(C)\), let \(\Fin(\tau)\) be the closed ideal
generated by the positive elements on which \(\tau\) is finite,
and let \(\Ker(\tau)\) be its zero ideal. The restriction of
\(\tau\) to \(\Fin(\tau)\) is densely finite, whereas
\(\tau(a)=\infty\) for \(a\in C_+\setminus\Fin(\tau)\).
A useful characterization is
\begin{equation}\label{st:finite-rank-cuts}
 a\in\Fin(\tau)
 \quad\Longleftrightarrow\quad
 d_\tau((a-\eps)_+)<\infty\quad\hbox{for every }\eps>0.
\end{equation}
To verify the forward direction, approximate \(a\) in norm by a
finite-trace positive element \(b\) of \(\Fin(\tau)\).
If \(\|a-b\|<\eps/2\), cut comparison gives
\[
 d_\tau((a-\eps)_+)
 \leq d_\tau((b-\eps/2)_+)
 \leq \frac{2}{\eps}\tau(b)<\infty.
\]
The finite-trace positives are dense in this ideal: finite sums of
terms \(x^*bx\), with \(\tau(b)<\infty\), still have finite trace
and generate its positive cone densely. Conversely, finite rank of
all positive cuts implies finite trace of those cuts, since
\(\tau((a-\eps)_+)\leq\|a\|d_\tau((a-\eps)_+)\).
Their norm limit is \(a\), proving \eqref{st:finite-rank-cuts}.

\begin{proposition}[Transport of ideals]\label{st:ideals}
The maps \(F,G\) identify the complete ideal lattices of \(P\) and
\(Q\), and these identifications agree. They identify the finite
and zero ideals of corresponding weights. The resulting
homeomorphism of primitive ideal spaces will be used to regard
\(\Prim(P)=\Prim(Q)=X\). For an open \(U\subseteq X\), write
\(P(U)\) and \(Q(U)\) for the corresponding ideals. Then
\[
 G(\Cu(P(U)))=\Cu(Q(U)),
\]
and
\[
 \Fin(\tau)=Q(U)\ \Longleftrightarrow\ \Fin(F\tau)=P(U),
 \qquad
 \Ker(\tau)=Q(U)\ \Longleftrightarrow\ \Ker(F\tau)=P(U).
\]
\end{proposition}

\begin{proof}
The additive idempotents of the trace cone are exactly the ideal
weights \(\tau_I\). Indeed \(\tau+\tau=\tau\) forces every
value to be zero or infinity, and the zero set is a closed ideal.
Moreover
\[
 I\subseteq J\quad\Longleftrightarrow\quad
 \tau_I+\tau_J=\tau_I.
\]
Thus \(F\) determines an order isomorphism of ideal lattices.
If \(F(\tau_I)=\tau_J\), then
\[
 [a]\in\Cu(J)
 \quad\Longleftrightarrow\quad
 d_{F\tau_I}(a)=0
 \quad\Longleftrightarrow\quad
 d_{\tau_I}(G[a])=0
 \quad\Longleftrightarrow\quad
 G[a]\in\Cu(I).
\]
This proves compatibility with \(G\).

The topology recovers the two other ideals through
\begin{equation}\label{st:scalar-limits}
 \lim_{r\downarrow0}r\tau=\tau_{\Fin(\tau)},
 \qquad
 \lim_{r\uparrow\infty}r\tau=\tau_{\Ker(\tau)}.
\end{equation}
These are \cite[Equation~(3.4)]{ERS}. They can also be checked
directly in \eqref{intro:trace-topology}: all positive cuts of an
element of \(\Fin(\tau)\) have finite trace by
\eqref{st:finite-rank-cuts}, whereas an element outside that ideal
has infinite trace; for the second limit use whether its trace is
zero or positive. Continuity and positive homogeneity of \(F\)
therefore preserve both limits.

An ideal-lattice isomorphism preserves prime ideals and their
hull--kernel topology. In a separable \(C^*\)-algebra prime ideals
are primitive \cite[Corollary~1]{DixmierPrime}, which gives the
asserted homeomorphism.
Finally, the zero algebra is distinguished by the cone: every nonzero
algebra has the nonzero weight \(\tau_{\{0\}}\). We may henceforth
assume that \(P,Q\ne0\).
\end{proof}

For a positive element \(c\in Q(U)\), we say that \(c\) is
\emph{compactly Cuntz supported in \(Q(U)\)} if
\begin{equation}\label{st:compact-controller}
 [c]\leq[(b-\delta)_+]
 \quad\hbox{for some }b\in Q(U)_+,\ \delta>0.
\end{equation}
Equivalently, \([c]\) is below an element compactly contained in a
Cuntz class from \(Q(U)\). This concerns the class of \(c\), not
compactness of a subset of a possibly non-Hausdorff space.

The next lemma permits one controller to be chosen before an
approximation sequence is constructed. Its finite multiplicity may
depend on the element being controlled. Later we will compensate for
that multiplicity by choosing a sufficiently late coordinate.

\begin{lemma}[Countable supports and common controllers]
\label{st:controllers}
The open-set lattice of \(X\) has a countable basis closed under
finite unions such that every open \(U\) is the union of basis opens
\(V\Subset U\). If \(V\Subset U\), there is a compactly Cuntz
supported \(b_{V,U}\in Q(U)_+\) with
\[
 Q(V)\subseteq\Ideal_Q(b_{V,U}).
\]
Every compactly Cuntz supported \(c\in Q(V)_+\) then satisfies
\begin{equation}\label{st:controller-multiple}
 [c]\leq N[b_{V,U}]
 \quad\hbox{for some finite }N\in\N.
\end{equation}
The element \(b_{V,U}\) is independent of \(c\); the integer
\(N\) need not be.
\end{lemma}

\begin{proof}
Choose a norm-dense sequence \((a_j)\) of positive contractions in
\(Q\). The supports of \((a_j-r)_+\), for positive rational
\(r\), form a countable open basis, after taking finite unions.
For example, to find such a support inside \(U\) and containing
\(x\in U\), choose \(a\in Q(U)_+\) nonzero in the quotient at
\(x\), and approximate it sufficiently closely by some \(a_j\).
A rational cut larger than the approximation error belongs to
\(Q(U)\), and a cut smaller than the quotient norm remains nonzero
at \(x\).
The support of \((a_j-r)_+\) is compactly contained, in the ideal
lattice, in the support of \((a_j-r/2)_+\). Indeed, if a directed
union of ideals contains the latter element, norm approximation
puts the more deeply cut element in one member of that union.
Choosing two rational cut levels inside the preceding construction
gives a basis open compactly contained in the prescribed \(U\).

Let \(h\) be strictly positive in \(Q(U)\). The ideals generated
by \((h-1/k)_+\) increase with dense union \(Q(U)\).
Compact containment gives \(k\) such that
\(Q(V)\subseteq\Ideal_Q((h-1/k)_+)\); take this cut as
\(b_{V,U}\). If \(c\) satisfies \eqref{st:compact-controller}
inside \(Q(V)\), then \([c]\ll[z]\) for some \(z\in Q(V)_+\).
The Cuntz-semigroup ideal generated by \([b_{V,U}]\) consists of
the elements below \(\sup_N N[b_{V,U}]\). Thus
\([z]\leq\sup_N N[b_{V,U}]\), and compact containment proves
\eqref{st:controller-multiple}.
\end{proof}

For the domains used below,
\[
 S_Y=C_0(Y)\otimes P,
 \qquad S_Y(U)=C_0(Y)\otimes P(U).
\]
Their primitive spaces are \(Y\times X\). For \(s\in S_Y\),
let \(U_s\) be the projection to \(X\) of the open support of
\(\Ideal_{S_Y}(s)\). Projection is open, and
\begin{equation}\label{st:smallest-label}
 s\in S_Y(U)\quad\Longleftrightarrow\quad U_s\subseteq U.
\end{equation}
Thus every element has a smallest coarse label. This is the
lower-semicontinuity condition on the domain ideal assignment used
in absorption. The fine labels on \(S_Y\), when used later, are
its actual open subsets of \(Y\times X\).

\subsection{Folding and rational averaging}

There are two matrix operations in the proof. Stabilization adds
internal matrix corners, whose ranks add with no normalization.
Convex averaging instead uses a separate UHF factor, whose normalized
trace supplies the convex coefficients. The following identifications
allow both operations to have range \(Q\) while retaining this
distinction and aligning any already chosen finite tests.

\begin{lemma}[Trace-preserving folding]\label{st:folding}
Finite matrix amplifications, including amplifications whose sizes
vary with the coordinate, may be folded into \(Q\) so as to preserve
all trace and rank calculations with unnormalized internal matrix
traces. The folding of a fixed first corner can be made arbitrarily
close to the identity on any prescribed finite subset of \(Q\),
by a multiplier-unitary adjustment.

If \(h_1,\ldots,h_l\) are \cpc maps with values in finite matrix
amplifications of \(Q\), and \(\lambda_i>0\) are rational with
\(\sum_i\lambda_i=1\), their normalized diagonal average is a
\cpc map \(h\) into \(Q\) satisfying, for \(s\geq0\),
\begin{align}
 \tau(h(s))&=\sum_i\lambda_i\tau(h_i(s)),
 &\tau(h(s)^2)&=\sum_i\lambda_i\tau(h_i(s)^2),
 \label{st:average-traces}\\
 d_\tau((h(s)-\eps)_+)
 &=\sum_i\lambda_i d_\tau((h_i(s)-\eps)_+).
 \label{st:average-ranks}
\end{align}
The traces on the right include the unnormalized internal matrices.
All equalities include zero and infinite values.
\end{lemma}

\begin{proof}
Use \eqref{st:W-properties} once to exhibit
\(Q=C\otimes\UHF\otimes\K\). For a matrix amplification choose
a Hilbert-space unitary
\(U_k:\C^k\otimes\ell^2\to\ell^2\) and fold by
\(\id_{C\otimes\UHF}\otimes\Ad(U_k)\).
This preserves the unnormalized matrix extension of every trace.
For a finite set \(E\subset Q\), choose a finite-rank projection
\(e\) in the last factor such that
\(\|a-eae\|\) is small for \(a\in E\).
A unitary of \(B(\ell^2)\) can send the folded first-corner copy
of \(e\ell^2\) back to \(e\ell^2\), agreeing with the specified
identification there. Conjugating the folding by this unitary fixes
the compressed finite tests. The two compression errors give the
desired approximation on \(E\).

Write \(\lambda_i=l_i/L\). In an additional UHF factor choose
orthogonal projections \(e_i\) with normalized traces
\(l_i/L\) and sum one. After folding the internal matrices, form
\[
 h(s)=\sum_i h_i(s)\otimes e_i.
\]
Orthogonality makes its norm the maximum of the block norms, and
functional calculus is blockwise for functions vanishing at zero.
This proves complete positivity, contractivity, and
\eqref{st:average-traces}--\eqref{st:average-ranks} before the final
UHF identification.

Choose that identification by interleaving tensor factors of a fixed
infinite tensor-product realization of \(\UHF\). Finite tensor
permutation unitaries identify its first-factor embedding with the
identity to any prescribed accuracy on finite tests. For a finite
matrix direction the trace identity on matrix units gives
\(\widetilde\tau(a\otimes e)=\operatorname{tr}(e)\tau(a)\),
where \(\tau(a)=\widetilde\tau(a\otimes1)\) and
\(\operatorname{tr}\) is normalized. Increasing finite matrix
algebras and positive cuts give this factorization for the UHF
direction, including extended weights. Thus the identification
preserves the displayed formulas even when \(Q\) is nonunital.
Multiplier conjugation also preserves every extended trace, by the
trace identity applied to \(a^{1/2}u\).
The choices are made separately at each coordinate; finite sets,
matrix sizes, and accuracies may therefore vary with that coordinate.
\end{proof}

We always distinguish the internal matrix directions from the
normalized candidate-averaging direction. In particular, an internal
sector sum is not divided by its number of sectors. All subsequent
padding and foldings use \Cref{st:folding} with the existing constant
and finite test elements included among the elements to be aligned.
